\section{本章小结}
\label{sec:fapprox-summary}

本章开头问：哪些简单函数能以任意精度逼近目标，加入表示预算后最佳误差又怎样随预算变化。对第一个问题，答案依赖测度、积分指数与误差范数。在任意测度空间中，截断与二进分层的简单函数在$L^p$（$1\leq p<\infty$）中稠密；在Lebesgue区域上，有理矩形与有理复系数给出可数稠密族，$L^p$可分，一般测度空间则要另核集合族条件。欧氏开集上，开紧逼近给出紧支撑连续函数的逼近，平移在$L^p$中连续，磨光误差由平移误差控制，于是光滑紧支撑函数也稠密，局部结论通过截断得到。这些结论在$L^\infty$中都失效：$L^\infty[0,1]$不可分，连续函数与跳跃函数的距离不小于跳跃的一半，尽管简单函数仍能按本质上确界逼近有界函数。

这些都是定性结论。为了比较效率，第~\ref{sec:fapprox-error}节用到候选集合的距离定义最佳逼近误差，并区分最佳误差与具体构造的误差、单个函数的误差与整类函数的最坏误差；后者在有限候选集合下正是覆盖数的问题。有限维子空间中最佳逼近总存在，Hilbert空间中由正交投影唯一给出，其他范数下可以不唯一，最佳选择的映射也不必是线性的。Fourier截断的平方误差等于被舍弃系数的平方和，一致误差由绝对尾和控制，两者要求的系数衰减不同；离散计算还带来与截断分开记账的混叠误差。连续频带截断给出带限函数中的最佳$L^2$逼近，但带限函数空间是无限维的，有限计算还需进一步离散。

后三节比较了三种候选集合。固定子空间对所有目标相同：多项式由次数控制，Bernstein构造给出连续模控制的收敛，附加光滑性带来更快的最佳速率；固定网格上的分段多项式由网格尺寸和每个单元内的局部正则性控制，拼接条件减少维数而不改变阶；张量积让误差随维数线性增长，系数个数却按$m^d$增长。最佳$N$项逼近固定正交基但按目标挑选支撑，在$L^2$中保留最大的$N$个系数最优，$\sigma_N(f)_2^2=\sum_{n>N}(c_n^*)^2$；结论依赖正交性与Parseval等式，不能直接搬到冗余字典或其他范数。自适应分段让分点随目标移动，对单调函数误差精确等于总变化量除以$2N$，对$\sqrt x$把速率从$N^{-1/2}$提高到$N^{-1}$；最佳分割的存在、误差界与找到分割的算法是三个不同的问题。

本章的误差估计都需要对目标的某种正则性作出假设：连续模、有界导数、系数衰减或总变化量。随后三章分别给出系统的刻画：第~\ref{chap:sobolev-spaces}章用弱导数与积分范数控制光滑性，第~\ref{chap:besov-spaces}章用不同尺度上的局部变化刻画非线性逼近的速率，第~\ref{chap:rkhs}章从离散点值观测出发研究插值与逼近。
